基于 CLIP 的零样本异常检测通常以固定的正常/异常文本原型解释目标类别，但固定原型难以描述单张测试图中的局部外观。直接从测试图构造视觉参照同样存在风险：仅依赖 CLIP 语义分数选择 patch 会继承其局部盲区，而直接把高频响应当作异常分数又会把规则纹理和物体边界误判为缺陷。本文提出一种\textbf{频谱可靠性引导的测试时语义原型校准}方法。我们在冻结的 CLIP patch 特征上施加 Haar 小波分解，以高频能量与低频结构边界的组合构造逐 patch 可靠性；随后将该可靠性作为语义证据的弱调制项，分别汇聚逐图正常与异常视觉原型，并以保守、非对称的方式校准固定文本原型。最终异常图仍由校准后的语义原型与 patch 特征的相似度产生，而非由频率图直接融合得到。该设计把小波从“异常评分器”重新定位为“逐图证据筛选器”，从而在保持 CLIP 语义空间可解释性的同时引入实例外观参照。本文给出面向五个数据集的主结果、直接融合负控、语义自参照对照、频谱设计消融、正常图稳定性和类别级分析方案。由于当前稿件用于冻结实验判据，表格中的数值为预期目标占位，正式版本将由真实评测结果替换。
